\chapter{Pre-processing e Costruzione del Dataset}
\label{ch:preprocessing}

\section{Descrizione del dataset}
\label{sec:dataset_desc}

Il fondamento empirico di questo lavoro di tesi è costituito da un dataset storico di natura azionaria, focalizzato sull'indice S\&P 500, assunto come proxy rappresentativa del mercato azionario globale e dell'economia statunitense. Al fine di garantire una base dati sufficientemente ampia per l'addestramento dei modelli di \emph{deep learning} e, contestualmente, catturare molteplici regimi di mercato (incluse fasi di espansione, recessione e shock sistemici), l'orizzonte temporale di riferimento è stato fissato tra l'anno 2000 e il 2020.

Per assicurare la continuità e la stabilità statistica del dataset, l'intera popolazione dei titoli è stata filtrata applicando un rigoroso requisito di sopravvivenza ininterrotta (\emph{survivorship}). Sono stati mantenuti esclusivamente gli asset che hanno registrato quotazioni giornaliere continue per tutto il ventennio di osservazione. Questo processo di epurazione ha prodotto un bacino finale di $N = 362$ titoli azionari. A differenza di campionamenti ridotti, il mantenimento dell'intera popolazione sopravvissuta garantisce una rappresentazione macrostrutturale completa del mercato, eliminando potenziali bias di selezione e preservando in modo organico l'eterogeneità settoriale GICS (\emph{Global Industry Classification Standard}). Tale eterogeneità è fondamentale affinché le matrici di correlazione empiriche manifestino la tipica struttura gerarchica e lo spettro degli autovalori aderente alla vera realtà economica.

I dati grezzi sono costituiti dai prezzi di chiusura giornalieri rettificati (\emph{adjusted closing prices}), i quali incorporano ex-ante gli effetti di operazioni societarie straordinarie, quali stacchi di dividendi, \emph{stock split} e raggruppamenti, prevenendo salti artificiali nelle serie storiche che inquinerebbero il calcolo delle covarianze.

\section{Calcolo dei rendimenti e stima delle matrici empiriche}
\label{sec:rendimenti_finestre}

A partire dalla matrice dei prezzi storici $\mathbf{P} \in \mathbb{R}^{N \times T_{\text{tot}}}$, dove $T_{\text{tot}}$ rappresenta il numero totale di giorni di contrattazione nel ventennio e $N=362$, il primo passaggio di trasformazione algebrica consiste nel calcolo dei log-rendimenti giornalieri continui. Il rendimento dell'asset $i$-esimo al tempo $t$ è definito come:
\begin{equation}
r_{i,t} = \ln \left( \frac{P_{i,t}}{P_{i,t-1}} \right)
\end{equation}

L'addestramento di modelli di \emph{machine learning} richiede un dataset composto da centinaia di campioni. Poiché la serie storica produce un'unica traiettoria temporale, si rende necessario implementare una procedura di estrazione a finestre mobili (\emph{rolling windows}).

Sia $T_w$ l'ampiezza temporale della finestra di osservazione, fissata operativamente a 724 giorni lavorativi (corrispondenti a quasi tre anni solari di contrattazioni). Definendo un parametro di passo (\emph{stride}) $S$, la finestra viene fatta scorrere lungo l'intera serie storica. Per ogni step $k$, si isola un sotto-campione di log-rendimenti $\mathbf{X}^{(k)} \in \mathbb{R}^{N \times T_w}$ e si calcola la corrispondente matrice di correlazione di Pearson empirica $\mathbf{C}^{(k)} \in \mathbb{R}^{N \times N}$, il cui generico elemento è:
\begin{equation}
C_{ij}^{(k)} = \frac{\text{Cov}(r_i, r_j)}{\sigma_i \sigma_j} \quad \text{con dati estratti in } t \in [k \cdot S, k \cdot S + T_w]
\end{equation}

Applicando un passo $S=10$ giorni (corrispondente esattamente a due settimane lavorative), si bilancia la necessità di generare un numero congruo di campioni per l'addestramento con l'esigenza di ridurre la ridondanza informativa del micro-rumore giornaliero. Questa estrazione tramite finestre ampie cattura in modo molto stabile e privo di rumore transitorio i veri \emph{shift} strutturali di medio-lungo periodo del mercato. In ogni finestra, il rapporto di dimensionamento temporale (rapporto di spettro) si attesta esattamente a $q = N/T_w = 362/724 = 0.5$. Tale valore, mantenendosi rigorosamente costante nel tempo, determina un livello di rumore statistico omogeneo e perfettamente bilanciato, il cui comportamento spettrale è analiticamente governabile dalle leggi della \emph{Random Matrix Theory} e dai limiti di Marchenko-Pastur discussi nei capitoli precedenti.
\section{Tecnica di campionamento: il Random Shuffle per la stazionarietà}
\label{sec:purging}

Nella strutturazione convenzionale di serie storiche finanziarie per il \emph{deep learning}, la prassi consolidata prevede una rigida divisione cronologica dei dati in \emph{Training Set}, \emph{Validation Set} e \emph{Test Set}, accompagnata da tecniche di \emph{purging} per gestire la sovrapposizione temporale causata dalle finestre mobili. Tuttavia, test preliminari condotti mediante partizionamento temporale sequenziale hanno evidenziato una dinamica di \emph{distribution shift} estrema: i regimi macroeconomici passati non fornivano informazioni topologiche sufficienti per permettere alle architetture profonde di generalizzare fuori campionario sui regimi futuri.

Per valutare il reale potere ricostruttivo e la capacità di compressione geometrica delle reti neurali, isolandole dall'impatto dei mutamenti macroeconomici radicali nel tempo, la pipeline sperimentale adotta una strategia di campionamento tramite mescolamento casuale (\emph{random shuffle}) sull'intero pool di matrici. 

Il processo di estrazione a finestre mobili, applicando il passo costante (\emph{stride}) $S=10$ giorni lavorativi lungo l'intera serie storica dei prezzi, produce un bacino complessivo di 412 matrici di correlazione empiriche. Queste matrici vengono interamente mischiate in modo stocastico prima di essere ripartite nei tre set operativi, portando alla seguente conformazione tensoriale finale:
\begin{itemize}
    \item \textbf{Training Set:} \emph{shape} $(288, 362, 362)$
    \item \textbf{Validation Set:} \emph{shape} $(82, 362, 362)$
    \item \textbf{Test Set:} \emph{shape} $(42, 362, 362)$
\end{itemize}

Questo approccio garantisce che la distribuzione statistica sottostante ai tre sottoinsiemi sia perfettamente omogenea e stazionaria: scenari di elevata volatilità e scenari di mercato stabili risultano equamente distribuiti sia nel set di addestramento che in quelli di validazione e test. Questa omogeneizzazione annulla l'impatto del \emph{distribution shift} e permette di testare rigorosamente la capacità dei modelli (PCA, Linear AE, Deep AE, VAE) di apprendere la complessa \emph{manifold} delle correlazioni a parità di condizioni statistiche.

\section{Estrazione e flattening delle matrici triangolari inferiori}
\label{sec:flattening}

L'ultimo passaggio della pipeline riguarda la conformazione dell'input per le reti neurali. Ogni matrice di correlazione empirica generata $\mathbf{C}^{(k)}$ ha dimensione $N \times N$, ovvero $362 \times 362$. Tuttavia, per le proprietà intrinseche dell'operatore di correlazione, la matrice è perfettamente simmetrica rispetto alla diagonale principale ($C_{ij} = C_{ji}$) e presenta esclusivamente valori unitari sulla diagonale stessa ($C_{ii} = 1$).

Fornire in pasto a un modello di \emph{deep learning} l'intera griglia $362 \times 362$ comporterebbe un grave spreco di risorse computazionali e introdurrebbe una ridondanza geometrica dannosa, costringendo il modello ad allocare pesi sinaptici per apprendere relazioni di simmetria identiche e valori costanti banalmente noti.

Per ottimizzare drasticamente lo spazio delle feature, si procede all'isolamento e all'estrazione della sola matrice triangolare inferiore (\emph{strictly lower triangular matrix}), scartando sia la diagonale principale sia la porzione superiore. Il numero $D$ di elementi unici e informativi estratti da una matrice $N \times N$ segue la formula delle combinazioni semplici:
\begin{equation}
D = \frac{N(N - 1)}{2}
\end{equation}

Essendo il paniere azionario composto da $N = 362$ titoli, ogni singola matrice viene compressa in un vettore di dimensione spaziale estremamente elevata:
\begin{equation}
D = \frac{362 \times 361}{2} = 65341 \text{ elementi}
\end{equation}

Attraverso un'operazione di \emph{flattening}, gli elementi della matrice triangolare inferiore vengono linearizzati in un vettore monodimensionale $\mathbf{x}^{(k)} \in \mathbb{R}^{65341}$. Questo vettore strutturato e formattato costituisce a tutti gli effetti l'input canonico $\mathbf{x}$ ad altissima dimensionalità che alimenterà i \emph{layer} di ingresso della PCA, dei modelli deterministici (Linear AE, Deep AE) e probabilistici (VAE) nei successivi step sperimentali.